\section{Exact data for one support, both phases}\label{sec:ppt3-app-coexistence}

This appendix contains the proofs of Theorem~\ref{thm:ppt3-phase-coexistence} and
Proposition~\ref{prop:ppt3-exposed-transition}, with the exact rational coefficient matrices and principal minors. The
notation, including $\Gamma_z$ of \eqref{eq:ppt3-gamma}, is that of Section~\ref{sec:ppt3-coexistence}.

In the basis $(u,v,w,e,o)$,
\begin{equation}\label{eq:ppt3-bs-same-support}
 X_{\rm BS}=\frac1{54375}
 \begin{pmatrix}
 1232&0&0&0&0\\
 0&30493&32&0&0\\
 0&32&1693&0&0\\
 0&0&0&3700&0\\
 0&0&0&0&925
 \end{pmatrix}.
\end{equation}

\begin{proof}[Proof of Theorem~\ref{thm:ppt3-phase-coexistence}]
The norms
$\|Bq_1(2)\|^2=1250$ and
$\|Bq_1(2i)\|^2=290$ give by direct expansion
\[
 \sigma_{\rm BS}
   =\tfrac13\bigl(\ket{111}\bra{111}+\Gamma_2+\Gamma_{2i}\bigr)
   =BX_{\rm BS}B^\dagger ,
\]
with $X_{\rm BS}$ as in \eqref{eq:ppt3-bs-same-support}; $\sigma_{\rm BS}$ is biseparable because each
$\Gamma_z$ is. The coefficient matrix
is positive definite and, after applying $B$, gives a trace-one state on $R$.
The state is invariant under cyclic permutations and global parity. For the unnormalised
matrix $54375\,\sigma_{\rm BS}^{T_A}$,
the leading principal minors of its even and odd
parity blocks, in the orders of
\eqref{eq:ppt3-minors}, are respectively
\[
\begin{array}{c|rrrr}
&\Delta_1&\Delta_2&\Delta_3&\Delta_4\\\hline
\mathrm{even}&1232&4558400&11751097500&13766553000000\\
\mathrm{odd}&2618&4698075&8087598750&83202897843750
\end{array}.
\]
All are strictly positive, so Sylvester's criterion
and cyclic invariance give
$\sigma_{\rm BS}^{T_c}\succ0$ for each cut.
The exact expansion and minors are checked
independently (Appendix~\ref{sec:ppt3-verifiers}).

The state $\rho_{10/3}$ has range $R$,
rank five and strictly positive partial transposes,
but is GME by \eqref{eq:ppt3-h5-value}.
For the interior assertion, let $\mu$ be a probability
measure of full support on the complex projective line $\mathbb{CP}^1$ (for example
normalised Fubini--Study measure) and
$M=\int q_1(z)q_1(z)^\dagger\,d\mu/\|Bq_1(z)\|^2$, the
coefficient matrix of a biseparable state on the
$A|BC$ product curve. If a Hermitian $H$ has
$q_1(z)^\dagger Hq_1(z)\ge0$ for all $z$ and
$\tr(HM)=0$, the nonnegative integrand vanishes on the full support of the measure, so the polynomial vanishes identically. The
five entries of $q_1(z)$ form a basis of the polynomials of degree at most four. Their products are therefore an
invertible change of basis from the independent monomials $\bar z^i z^j$, $0\le i,j\le4$, so $H=0$. Hence $M$ is
interior to the cone of supported $A|BC$-biseparable
coefficient matrices. (A measure on the unit circle
would not do: $H_5$ vanishes there.)
Mixing a sufficiently small amount
of such a moment matrix into $\sigma_{\rm BS}$
produces a biseparable state interior to the
supported biseparable cone, while retaining
strict PPT and rank five. GME is open and $\rho_{10/3}^{T_c}\succ0$,
so a relative neighbourhood of $\rho_{10/3}$
consists of fully-PPT GME states. Both regions therefore have relative interior in the same PPT spectrahedron.
\end{proof}

In the basis $(u,v,w,e,o)$,
\begin{equation}\label{eq:ppt3-xdelta}
 X_\partial=\frac1{3360}
 \begin{pmatrix}
 293&0&0&0&0\\
 0&577&-105&0&0\\
 0&-105&273&0&0\\
 0&0&0&188&0\\
 0&0&0&0&76
 \end{pmatrix}.
\end{equation}

\begin{proof}[Proof of Proposition~\ref{prop:ppt3-exposed-transition}]
For the rays in \eqref{eq:ppt3-exposed-state},
$|z|=1$ at $z=\pm1,\pm i$, while at
$z=\pm i\sqrt5$ one has $r=5$ and
$\operatorname{Re}z=0$. The factorisation
\eqref{eq:ppt3-h5} therefore vanishes on every
summand, proving $\tr(H_5X_\partial)=0$.
Their squared physical norms
$\|Bq_1(1)\|^2=56$, $\|Bq_1(i)\|^2=32$
and $\|Bq_1(i\sqrt5)\|^2=576$ give the
coefficient matrix \eqref{eq:ppt3-xdelta} upon
cyclic and sign averaging. It has
$\tr(BX_\partial B^\dagger)=1$ and
$X_\partial\succ I_5/100$ (a diagonal and
$2\times2$ block check).

For a quantitative PPT bound, the leading
principal minors of the even and odd
$4\times4$ parity blocks of
$5\,(3360\,\sigma_\partial^{T_A})-56I_8$,
in the orders of \eqref{eq:ppt3-minors}, are
\[
\begin{array}{c|rrrr}
 &\Delta_1&\Delta_2&\Delta_3&\Delta_4\\\hline
\mathrm{even}&1409&1163516&654702384&109183407616\\
\mathrm{odd}&1689&1363796&97132144&7881970176
\end{array}.
\]
They are positive. Hence
$\sigma_\partial^{T_A}\succ I/300$ by
Sylvester's criterion, and cyclic symmetry
gives the other cuts.

Let $\Delta=\sigma_{\rm BS}-\sigma_\partial$
and $X_\Delta=X_{\rm BS}-X_\partial$.
Direct rational arithmetic gives
\[
 \|\Delta\|_2^2=
 \frac{575589900887}{2318006250000}<\frac14 .
\]
Partial transposition preserves the
Hilbert--Schmidt norm, so
$\|\Delta^{T_c}\|_\infty<1/2$.
Likewise $B^\dagger B=\operatorname{diag}(4,1,3,3,3)
\succeq I_5$ implies
$\|X_\Delta\|_\infty\le\|X_\Delta\|_2
\le\|\Delta\|_2<1/2$.
For $-1/200\le t\le0$, it follows that
\[
 \sigma(t)^{T_c}\succ
   \left(\frac1{300}-\frac1{400}\right)I
      =\frac1{1200}I,\qquad
 X_\partial+tX_\Delta\succ
   \left(\frac1{100}-\frac1{400}\right)I_5 .
\]
For $0\le t\le1$ strict PPT and rank five
follow by convexity from the corresponding
properties of the two endpoints. Biseparability in that range follows from the explicit decompositions of the endpoints.
Finally, for $t<0$ the supported witness has
\[
 \tr\!\bigl(H_5(X_\partial+tX_\Delta)\bigr)
 =t\,\frac{25082}{54375}<0,
\]
which excludes biseparable mixtures.
The arithmetic is checked independently (Appendix~\ref{sec:ppt3-verifiers}).
\end{proof}
